\documentclass[conference]{IEEEtran}

\usepackage{cite}
\usepackage{url}

\begin{document}
	
	\title{Digital Image Processing Summary for License Plate Recognition and Vehicle Counting}
	
	\author{
		\IEEEauthorblockN{Muhammad Fathan Romiza }
		\IEEEauthorblockA{
			\textit{Digital Image Processing Course} \\
			\textit{Gadjah Mada University} \\
			Yogyakarta, Indonesia \\
			muhammadfathanromiza560518@mail.ugm.ac.id
		}
	}
	
	\maketitle
	
	\begin{abstract}
		Digital image processing can be used to extract information from
		images and videos for transportation applications. This report
		summarizes two cases, license plate recognition and vehicle counting,
		and identifies image features suitable for each task. In license plate
		recognition, vertical and horizontal edge densities help locate
		candidate plate regions, while connected component analysis supports
		character segmentation. In vehicle counting, foreground regions,
		contours, and object positions help identify vehicles and determine
		when they pass a virtual detection line or area. The discussion draws
		on five international journal articles, with two studies examined as
		the main cases. The comparison shows that feature selection depends on
		the task and the information required from the input image or video.
	\end{abstract}
	
	\begin{IEEEkeywords}
		Digital image processing, license plate recognition, vehicle counting,
		edge detection, image segmentation
	\end{IEEEkeywords}
	
	
	
	\section{Introduction}
	
	Digital image processing uses computational methods to improve images
	and extract information from them. In transportation, images and
	videos can provide information about vehicles and traffic conditions.
	Two applications are automatic license plate recognition and vehicle
	counting.
	
	License plate recognition aims to locate a vehicle's plate and
	identify its characters. The task involves detecting the plate region
	within an image, separating the characters, and recognizing each
	character. Changes in illumination, image quality, plate size, and
	background can affect the process \cite{arafat2020}.
	
	Vehicle counting aims to determine the number of vehicles passing
	through a road section. A video-based system can separate moving
	vehicles from the background and count them as they pass a predefined
	line or area. The task becomes more difficult when vehicles are close
	together, shadows are present, or lighting changes \cite{liu2017}.
	
	This report summarizes the two applications using international
	journal articles. It discusses the processing stages, identifies
	suitable image features, and explains how those features contribute to
	each task.

	
	\section{Background}
	
	\subsection{Image Features}
	
	An image feature is a measurable characteristic that describes part
	of an image. Features can represent properties such as intensity,
	color, texture, shape, size, or position. These properties can help a
	system locate an object, distinguish it from its surroundings, or
	compare it with another object.
	
	The appropriate features depend on the task. Edge information can help
	locate a region containing characters, while an object's position can
	help determine whether it has crossed a counting line.
	
	\subsection{Image Processing Techniques}
	
	Several image processing techniques are relevant to the selected cases.
	
	\textit{Grayscale conversion} represents an image using intensity
	values. It can simplify subsequent operations when color information is
	unnecessary for a particular processing stage.
	
	\textit{Edge detection} identifies changes in image intensity. These
	changes can indicate object boundaries or transitions between
	characters and their background.
	
	\textit{Segmentation} divides an image into regions of interest. In
	license plate recognition, segmentation can separate individual
	characters. In vehicle counting, it can help distinguish moving
	foreground regions from the background.
	
	\textit{Morphological processing} modifies image regions based on their
	shape and pixel relationships. Operations such as dilation can prepare
	an image for later detection or segmentation.
	
	\textit{Connected component analysis} groups neighboring foreground
	pixels into regions. These regions can then be examined as possible
	objects or characters.
	

	
	\section{Case Studies}
	
	\subsection{License Plate Recognition}
	
	\subsubsection{Problem and Objective}
	
	License plate recognition extracts a plate from an image of a vehicle
	and converts its characters into a sequence of letters and numbers.
	
	Arafat \textit{et al.} \cite{arafat2020} proposed a framework that
	combines edge-based detection and connected component analysis to
	locate plates and segment their characters. Their framework considers
	plates of different sizes and shapes from cars, buses, trucks, and
	motorcycles.
	
	\subsubsection{Image Processing Stages}
	
	The framework in \cite{arafat2020} consists of six main stages:
	
	\begin{enumerate}
		\item \textit{Preprocessing:} The input image is converted to
		grayscale and processed using dilation-based operations.
		
		\item \textit{Edge detection:} Vertical and horizontal edges are
		detected using kernel matrices.
		
		\item \textit{Energy mapping:} Edge information is used to produce
		vertical and horizontal energy maps. Gaussian smoothing helps
		identify candidate plate regions.
		
		\item \textit{Plate localization:} A candidate plate region is
		located and extracted.
		
		\item \textit{Character segmentation:} The extracted plate is
		converted into a binary image. Connected component analysis and
		blob assessment separate potential characters.
		
		\item \textit{Character recognition:} The segmented characters are
		resized and compared with stored character templates using a
		correlation-based matching method.
	\end{enumerate}
	
	The processing flow can be summarized as follows:
	
	\begin{quote}
		Vehicle image $\rightarrow$ preprocessing $\rightarrow$ edge detection
		and energy mapping $\rightarrow$ plate localization $\rightarrow$
		character segmentation $\rightarrow$ character recognition.
	\end{quote}
	
	\subsubsection{Suitable Image Features}
	
	The main features used to locate the plate are \textit{vertical and
		horizontal edge densities}. A license plate contains characters and
	boundaries that create intensity changes in the image. Measuring these
	changes in both directions helps the system identify regions that may
	contain a plate \cite{arafat2020}.
	
	After plate localization, the connected structure of foreground pixels
	becomes important. Connected component analysis groups neighboring
	pixels into candidate character regions. The system then compares the
	segmented character images with stored templates to determine their
	identities \cite{arafat2020}.
	
	Another approach to plate localization uses mathematical morphology to
	emphasize structures in a grayscale image. Yepez and Ko
	\cite{yepez2018} applied a top-hat transformation to help distinguish
	potential plate regions from their surroundings. Their method shows how
	the size and shape of a structuring element can be used to identify
	characteristics relevant to plate localization.
	
	Kim \cite{kim2019} combined a region-based convolutional neural network
	with image processing techniques for vehicle and license plate
	extraction. The method uses Sobel edge filtering and morphological
	operations during plate-region extraction. This study illustrates the
	use of edge information and morphology alongside a learned vehicle
	detection method.
	
	\subsubsection{Reported Results and Limitations}
	
	Using 400 images of Malaysian license plates, Arafat \textit{et al.}
	reported a 96.5\% plate localization rate, a 95.6\% character
	segmentation rate, and a 94.4\% character recognition rate. These
	percentages describe different stages of the framework
	\cite{arafat2020}.
	
	The authors reported errors in images affected by overexposure,
	darkness, blur, dirty plates, and complex backgrounds. Similar
	character shapes also caused recognition errors, such as confusing the
	number ``0'' with the letter ``O'' \cite{arafat2020}.
	
	
	\subsection{Vehicle Counting}
	
	\subsubsection{Problem and Objective}
	
	Vehicle counting estimates how many vehicles pass through a monitored
	road section. Liu \textit{et al.} \cite{liu2017} proposed a video-based
	system for urban roads that detects vehicles and adapts its counting
	method to traffic conditions.
	
	The system uses a virtual detection line when traffic moves freely and
	a virtual loop, or predefined image region, when traffic is congested.
	This choice addresses differences in vehicle movement and spacing under
	the two conditions \cite{liu2017}.
	
	\subsubsection{Image Processing Stages}
	
	The system in \cite{liu2017} processes frames from a road video using
	the following stages:
	
	\begin{enumerate}
		\item \textit{Background initialization and updating:} The system
		builds a representation of the road background and updates it as
		conditions change.
		
		\item \textit{Foreground extraction:} Background subtraction helps
		identify regions that may correspond to moving vehicles.
		
		\item \textit{Image segmentation and refinement:} Filtering and
		segmentation are applied, together with shadow suppression in
		daytime or lamplight suppression at night.
		
		\item \textit{Contour extraction:} The system extracts and fills
		contours to obtain vehicle regions.
		
		\item \textit{Vehicle counting:} The system selects a detection line
		or virtual loop according to the traffic state and updates the
		vehicle count.
	\end{enumerate}
	
	The processing flow can be summarized as follows:
	
	\begin{quote}
		Road video $\rightarrow$ background subtraction $\rightarrow$ vehicle
		segmentation $\rightarrow$ contour extraction $\rightarrow$ virtual
		line or loop $\rightarrow$ vehicle count.
	\end{quote}
	
	\subsubsection{Suitable Image Features}
	
	The system uses \textit{foreground regions and contours} to represent
	detected vehicles. After a vehicle region is obtained, its
	\textit{position} relative to a virtual detection line helps determine
	whether it has passed through the monitored location.
	
	The system also uses region size information in the virtual loop. The
	ratio of foreground pixels within the loop and the average width of the
	detected object help determine whether the loop is occupied by a
	vehicle \cite{liu2017}.
	
	These features suit the counting task because the system needs to
	observe the presence and movement of vehicle regions over time. For the
	line-based method, a count is updated as a vehicle occupies and then
	leaves the detection line \cite{liu2017}.
	
	Rosas-Arias \textit{et al.} \cite{rosas2019} proposed another
	video-based vehicle counting method using Incremental Principal
	Component Analysis to detect changes between video frames. The
	resulting image is thresholded and refined to identify moving regions.
	Vehicles are then counted using a virtual detection line in each lane.
	This approach also uses the presence of foreground pixels along a line
	as information for counting.
	
	\subsubsection{Reported Results and Limitations}
	
	Liu \textit{et al.} developed and tested a prototype on an urban road.
	They reported real-time counting accuracy above 99\% in most of their
	field scenarios \cite{liu2017}. This result describes the conditions
	evaluated in that study.
	
	The authors noted that nearby vehicles in congested traffic could be
	detected as a single object. They also stated that the system was not
	tested under severe conditions such as heavy fog, heavy rain, or snow
	\cite{liu2017}.
	
	

	\section{Comparison and Discussion}
	
	The two main cases use digital image processing to extract information
	about vehicles. Each task requires different information.
	Table~\ref{tab:comparison} summarizes the main differences between the
	two applications.
	
	\begin{table}[htbp]
		\caption{Comparison of the Two Main Applications}
		\label{tab:comparison}
		\centering
		\footnotesize
		\begin{tabular}{|p{0.18\columnwidth}|p{0.34\columnwidth}|p{0.34\columnwidth}|}
			\hline
			\textbf{Aspect} & \textbf{License Plate Recognition} &
			\textbf{Vehicle Counting} \\
			\hline
			Input & Vehicle image & Sequence of road video frames \\
			\hline
			Main object & License plate and characters & Vehicle regions \\
			\hline
			Main features & Edge density and connected character regions &
			Foreground region, contour, position, and size \\
			\hline
			Techniques & Edge detection, morphology, connected component analysis,
			and template matching & Background subtraction, segmentation, contour
			extraction, and virtual line or loop \\
			\hline
			Output & Recognized plate characters & Number of vehicles \\
			\hline
		\end{tabular}
	\end{table}
	
	In license plate recognition, edge features help locate a region
	containing closely arranged characters. Connected pixels and template
	similarity are then used during character segmentation and recognition.
	The processing stages move from a vehicle image toward individual
	characters \cite{arafat2020}.
	
	In vehicle counting, the system examines vehicle regions across video
	frames. Position and region size are suitable because the output
	depends on whether a vehicle passes a monitored line or occupies a
	virtual loop \cite{liu2017}.
	
	The additional studies show that methods can vary within the same
	application. Plate localization can use morphological features
	\cite{yepez2018} or combine a neural network with edge and
	morphological processing \cite{kim2019}. Vehicle counting can use an
	adaptive background model \cite{liu2017} or Incremental Principal
	Component Analysis for motion detection \cite{rosas2019}.
	
	The studies also show that preprocessing and segmentation influence
	later stages. Poor separation of a plate or character can affect
	recognition, while poor separation of nearby vehicles can affect
	counting \cite{arafat2020,liu2017}.
	
	
	
	\section{Conclusion}
	
	This report discussed license plate recognition and vehicle counting as
	two applications of digital image processing in transportation.
	
	For license plate recognition, vertical and horizontal edge densities
	support plate localization, while connected component analysis and
	template matching support character segmentation and recognition. For
	vehicle counting, foreground regions, contours, position, and size
	support the detection and counting of vehicles in video.
	
	The two cases show that image features should be selected according to
	the information required by the application. Their performance also
	depends on image conditions and the quality of the regions produced
	during processing.
	

	
	\begin{thebibliography}{00}
		
		\bibitem{arafat2020}
		M. Y. Arafat, A. S. M. Khairuddin, and R. Paramesran,
		``Connected component analysis integrated edge based technique for
		automatic vehicular license plate recognition framework,''
		\textit{IET Intelligent Transport Systems}, vol. 14, no. 7,
		pp. 712--723, 2020, doi: \url{https://doi.org/10.1049/iet-its.2019.0006}.
		
		\bibitem{liu2017}
		F. Liu, Z. Zeng, and R. Jiang,
		``A video-based real-time adaptive vehicle-counting system for urban roads,''
		\textit{PLOS ONE}, vol. 12, no. 11, Art. no. e0186098, 2017,
		doi: \url{https://doi.org/10.1371/journal.pone.0186098}.
		
		\bibitem{yepez2018}
		J. Yepez and S.-B. Ko,
		``Improved license plate localisation algorithm based on morphological operations,''
		\textit{IET Intelligent Transport Systems}, vol. 12, no. 6,
		pp. 542--549, 2018, doi: \url{https://doi.org/10.1049/iet-its.2017.0224}.
		
		\bibitem{kim2019}
		J. Kim,
		``Automatic vehicle license plate extraction using region-based
		convolutional neural networks and morphological operations,''
		\textit{Symmetry}, vol. 11, no. 7, Art. no. 882, 2019,
		doi: \url{https://doi.org/10.3390/sym11070882}.
		
		\bibitem{rosas2019}
		L. Rosas-Arias \textit{et al.},
		``Vehicle counting in video sequences: An incremental subspace learning approach,''
		\textit{Sensors}, vol. 19, no. 13, Art. no. 2848, 2019,
		doi: \url{https://doi.org/10.3390/s19132848}.
		
	\end{thebibliography}
	
\end{document}